% [YOURS] — write last, in your own voice, after the results are on the page.
% The skeleton below is the fact set, not the prose. Your July timeline set the
% standard: "a human-rewritten abstract ... no AI-scaffold phrasing."
\begin{abstract}
Pretrained EEG foundation models are increasingly proposed as general-purpose
encoders for brain--computer interfaces, but published audits disagree about
whether their pretraining transfers. We audit two such models, LaBraM and
CBraMod, on motor imagery under a \emph{validation-locked} protocol: every
tunable choice --- preprocessing, architecture, stopping epoch, temperature,
and which method is reported --- is made on training-session data alone and
held-out labels are read once, a constraint on model \emph{selection} rather
than on model parameters, enforced by the result format rather than by author
discipline. On four-class
BCI~IV-2a, every supervised comparator we test outperforms every
foundation-model configuration, including the strongest fine-tuning recipe
selected on validation. We then run the matched-input control that such
comparisons require, retraining each supervised comparator on the identical
broadband arrays the foundation models consume, and report the paper's main
methodological result: the measured preprocessing term differs in sign across
comparator architectures, helping one by 0.078 accuracy while costing another
0.088, though no individual term reaches significance after correction at
$n=9$. The term is therefore not reliably estimable from any single
architecture, which is how such controls are ordinarily run. The deficit is
task-dependent: on two-class BNCI2014\_004 we could not detect a
difference between a fine-tuned CBraMod and the supervised arms. Calibration is the one dimension on
which the foundation models hold up; after a single validation-fitted
temperature their held-out calibration error lands in the supervised range
despite far lower accuracy.
\end{abstract}

\begin{IEEEkeywords}
EEG, brain--computer interface, foundation models, motor imagery, model
evaluation, calibration
\end{IEEEkeywords}
